\chapter{Discussion}

\section{Interpretation of the Results}

The results indicate that fine-tuning has a positive effect on both similarity-based
and execution-based evaluation metrics. For both StarCoder2 and Qwen, the
fine-tuned variants achieved higher BLEU scores, more accepted solutions, and
fewer compilation errors than the corresponding base models. This suggests that
fine-tuning enables the models to better adapt to the specific characteristics of
the Java-to-C++ translation task.

The explanation for this improvement is that the fine-tuning process exposes the models to a
large number of parallel Java and C++ examples from the XLCoST dataset. As a
result, the models can learn common translation patterns, language-specific
syntax, and frequently occurring structures. Previous studies have similarly
shown that fine-tuning on code-specific datasets improves both syntactic
correctness and functional performance in code-related tasks \cite{peft_code, hu2021lora, zhu2022xlcost}.

The Qwen models consistently outperformed StarCoder2 in both BLEU score and
execution-based evaluation. One likely reason is that Qwen2.5-Coder has been
trained on a larger and more diverse code corpus and is more strongly optimized
for programming-related tasks. It also benefits from instruction tuning, which
helps it follow translation prompts more reliably and produce cleaner outputs.


Compilation errors remained the most common failure category for all models.
Across the evaluated samples, a substantial portion of incorrect outputs failed
during the compilation stage. These errors included syntax issues, missing
imports, incorrect type usage, and incompatibilities between source and target
language constructs. One reason for this is that source code translation requires
not only understanding the source program, but also producing fully valid code
in the target language. Small mistakes in punctuation, memory handling, or
language-specific constructs can cause the entire translation to fail.




\subsection{Relationship Between BLEU and Execution-Based Evaluation}

To examine the relationship between similarity-based and execution-based
evaluation more precisely, a correlation analysis was performed between BLEU
score and two binary execution-based outcomes: accepted solution and compilation
success. The accepted solution variable was set to 1 when the generated program
passed all tests, and 0 otherwise. Compilation success was set to 1 when the
generated program did not result in a compilation error, and 0 otherwise.

Both Pearson and Spearman correlations were computed. Pearson correlation was
used to measure the linear association between BLEU and the execution-based
outcomes, while Spearman correlation was used to measure whether higher BLEU
scores generally correspond to better-ranked execution outcomes. The results are
shown in Table~\ref{tab:bleu_execution_correlation}.

\begin{table}[b]
\centering
\caption{Correlation between BLEU score and execution-based outcomes}
\begin{tabular}{lcccc}
\hline
Execution-based outcome & Pearson $r$ & Pearson $p$ & Spearman $\rho$ & Spearman $p$ \\
\hline
Accepted solution & 0.510 & $1.85 \times 10^{-101}$ & 0.368 & $4.46 \times 10^{-50}$ \\
Compilation success & 0.372 & $4.81 \times 10^{-51}$ & 0.509 & $6.12 \times 10^{-101}$ \\
\hline
\end{tabular}
\label{tab:bleu_execution_correlation}
\end{table}

The correlation results show a statistically significant positive relationship
between BLEU score and both execution-based outcomes. In both cases, the
$p$-values are far below 0.001, which means that the null hypothesis of no
correlation can be rejected. Therefore, higher BLEU scores are not completely
independent from execution-based success: outputs with higher BLEU scores are
generally more likely to compile successfully and to be accepted.

However, the strength of the correlations also shows that BLEU is not a fully
reliable proxy for functional correctness. The correlation between BLEU and
accepted solutions is moderate according to Pearson correlation ($r = 0.510$),
but weaker according to Spearman correlation ($\rho = 0.368$). This suggests that
although BLEU has some relationship with final correctness, the ranking of
solutions by BLEU does not consistently match the ranking implied by accepted
execution results. In practical terms, a higher BLEU score increases the
likelihood of a correct translation, but it does not guarantee that the generated
program passes the test cases.

The relationship between BLEU and compilation success shows a slightly different
pattern. The Pearson correlation is lower ($r = 0.372$), while the Spearman
correlation is stronger ($\rho = 0.509$). This indicates a moderate monotonic
relationship: translations with higher BLEU scores are generally more likely to
compile, but the relationship is still far from deterministic. This is consistent
with the observation that fine-tuned models produced fewer compilation errors,
but still generated many wrong answers. In other words, fine-tuning often helped
the models produce syntactically more valid C++ code, but syntactic validity did
not always imply semantic correctness.

Therefore, the results support the use of BLEU only as an auxiliary metric.
Execution-based evaluation remains necessary because it directly checks whether
the generated C++ program compiles and produces the expected output. A complete
evaluation of code translation systems should therefore combine similarity-based
metrics with execution-based testing, rather than relying on BLEU as the sole
indicator of translation quality.

